\answer{ex:04:1}{(a)~$\approx1.1\cdot10^{11}$~bit/J。(b)~$\log_2\sqrt{10}\approx1.7$ビット対$81$ビットで、$\approx50$分の1。}
\answer{ex:04:2}{$r^*=(P_0/E_{\text{spk}})\ln\bigl(1/(r^*\Delta t)\bigr)$、$P_0/E_{\text{spk}}=1.16\,\text{s}^{-1}$より、$r^*\approx6$スパイク/s、$7.4\cdot10^{10}$~bit/J。}
\answer{ex:04:3}{$\mathcal I=\frac{N}{1+(N-1)c}=9.9$（極限は$10$）対$\mathcal I=\frac{N}{1-c}=1111$で、$\approx110$倍。相関が有害なのは、信号に似ているときだけである。}
\answer{ex:04:4}{$\theta\approx68.7$と$19.9$（$w$を単位として）、$m=19$と$m=10$。速い受け手は、平均入力が5分の1なのに、同時入力がほぼ半分で足りる。}
\answer{ex:04:5}{$U\tau_{\text{rec}}\ge0.725\,\text{s}$かつ$\tau_{\text{rec}}(1-2U)\le0.05\,\text{s}$で、これには$U\ge0.483$が必要。最小の$\tau_{\text{rec}}=0.725\,\text{s}$は$U=1$のとき（$U=0.5$なら$1.45\,\text{s}$）。}
\answer{ex:04:6}{(a)~$\theta\,e/(e-1)\approx1.58\,\theta$で、$\tau$と$\theta$によらない。(b)~スパイク$14$個。}
\answer{ex:04:7}{(a)~単語あたり$1.62$ビット：スパイク数が$0.77$、その配置が$0.84$。(b)~$1.13$ビットと$1.59$ビット。$30$個の単語からの推定はほぼ3分の1を失う。}
